\section{Related Work}
\label{sec:related}

\subsection{Single-Cell Representation and Clustering}

Standard single-cell workflows project expression, construct a cell graph, and
apply community detection such as Leiden~\cite{wolf2018scanpy,traag2019leiden}.
They are scalable and interpretable, but their output depends on preprocessing,
neighborhood size, resolution, and initial geometry~\cite{kiselev2019challenges}.
Deep models learn this geometry. scVI uses a probabilistic latent
variable model~\cite{lopez2018scvi}; scDeepCluster and scDCC couple nonlinear
representation learning to clustering structure~\cite{tian2019scdeepcluster,tian2021scdcc}.
Graph-aware methods such as scGNN and scDSC propagate structural information
through neural objectives~\cite{wang2021scgnn,gan2022scdsc}, while scNAME and
AttentionAE-sc add neighborhood contrast, mask estimation, or attention
mechanisms~\cite{wan2022scname,li2023attentionae}.  These studies demonstrate
the value of intercellular context but also make success depend on an estimated
graph or pseudo-label structure.

Masked autoencoding avoids requiring a clustering target during representation
learning.  Inspired by masked modeling in vision~\cite{he2022mae}, scMAE
replaces gene values, predicts the realized mask, and reconstructs the original
expression vector~\cite{fang2024scmae}.  scVICAR retains this backbone and its
clean-cell objective, but uses the cell graph to define an auxiliary corruption
distribution from which the original anchor is recovered.

Contrastive single-cell methods provide a complementary route to structured
representations. CLEAR perturbs expression profiles to define positive views,
whereas scNAME combines mask estimation with neighborhood contrast and a memory
bank~\cite{wan2022scname}. More recent scRCL models heterogeneous cell views and
cell--gene associations through several contrastive objectives
~\cite{peng2024scrcl}. These methods optimize embedding agreement or separation.
scVICAR instead reconstructs the unmodified anchor and uses no contrastive loss
in its final configuration, allowing the effect of local corruption to be
tested without an additional representation-separation objective.

\subsection{Neighborhood Learning and Graph Risk}

Cell graphs are useful because nearby cells often share aspects of type, state,
lineage, or function.  Yet sparsity, class imbalance, batch effects, and the
projection used to build the graph can create false and asymmetric edges.
Repeated propagation may also contract informative graph-frequency components,
blurring boundaries in a manner related to oversmoothing~\cite{li2018laplacian}.
Rare cells are particularly exposed because they have fewer same-population
neighbors.

scVICAR uses the graph to construct training examples and extracts final
embeddings directly from clean profiles. Edge affinity determines each
neighbor's contribution, and a cell-specific coefficient controls the amount
of mixing. This separation supports direct tests of edge weighting, cell-level
mixing, and graph contamination. We use the term topology-informed affinity
for this uncalibrated analytic score.

\subsection{Vicinal Learning and Anchor Recovery}

Mixup regularizes a predictor using convex combinations of observations and
targets~\cite{zhang2018mixup}; graph variants extend interpolation to structured
data~\cite{wang2021graphmixup}. Adjacent cells may still occupy distinct
biological states. scVICAR uses their local convex mixture as a corrupted input
and retains the anchor as target. The model learns to recover the anchor after
a controlled displacement.

The graph-vicinal displacement is aligned with observed cell--cell variation
and bounded by local geometry.  Under additional smoothness and small-residual
conditions, its auxiliary risk admits a directional-curvature interpretation.
Random mixing, alternative aggregation estimators, and controlled bad-edge
tests measure the consequences of inaccurate graph directions.

\subsection{Evaluation Beyond a Single Clustering Score}

Single-cell benchmarking shows that method rankings depend on datasets,
preprocessing, and evaluation protocol~\cite{xu2026scclubench}.  Comparisons are
especially difficult when feature counts, capacity, training budget, or
clustering oracles differ.  We therefore use a matched backbone, disclose
known-$K$ KMeans, retain fixed-resolution Leiden as secondary, and prohibit a
label-guided resolution sweep.

We use multiple seeds to quantify algorithmic variability and average them
within each dataset before confirmatory inference. Marker recovery,
marker-overlap annotation, rare-class recall, and low-label probing evaluate
properties that clustering agreement alone leaves unresolved. Each task is
reported separately.

Biological validation in single-cell clustering commonly combines marker-gene
specificity, differential-expression patterns, enrichment analysis, rare-cell
recovery, and trajectory or spatial-coherence analyses. The appropriate task
depends on the biological structure represented in each dataset. We therefore
separate marker coherence from predictive separability and avoid treating a
high clustering score as evidence that every downstream biological property
has improved.
